\chapter{Solutions: Midterm Exam, November 8, 2017}
\label{s:physics-2017-11-midterm}
The density of states $D(E)$ below counts states per unit energy in the entire sample of volume $V$. A factor $g_s$ explicitly specifies spin counting. We use $g_s=2$ for the unpolarized free-electron gas.

\section*{1. A change of primitive basis}
Let $T=(\vct t_1\ \vct t_2\ \vct t_3)$ and $S=(\vct s_1\ \vct s_2\ \vct s_3)$. The stated relations give $S=TA$, where
\[
A=\begin{pmatrix}2&-1&1\\-1&1&\alpha\\2&\alpha&1\end{pmatrix},
\qquad
\det A=-(2\alpha+1)(\alpha+1)=-2\alpha^2-3\alpha-1.
\]
Two primitive bases generate the same lattice precisely when their change-of-basis matrix is \emph{integer unimodular}: $A$ has integer entries and $\det A=\pm1$. Necessity follows because both sets of basis vectors must be integer combinations of the other set. Sufficiency follows from the adjugate formula: an integer matrix with determinant $\pm1$ has an integer inverse.

The determinant equation $\det A=+1$ gives $2\alpha^2+3\alpha+2=0$, whose discriminant is $-7$. The equation $\det A=-1$ gives
\[
\alpha(2\alpha+3)=0,
\qquad \alpha=0\ \text{or}\ -\frac32.
\]
The second root gives noninteger coordinates for $\vct s_2$ and $\vct s_3$ in the original primitive basis; equal cell volumes alone do not make it a basis of the same lattice. Therefore
\[
\boxed{\alpha=0}.
\]
As a direct sufficiency check, at $\alpha=0$ the inverse relations are
\[
\vct t_1=-\vct s_1-\vct s_2+2\vct s_3,\qquad
\vct t_2=-\vct s_1+2\vct s_3,\qquad
\vct t_3=\vct s_1+\vct s_2-\vct s_3.
\]
They express every original lattice point as an integer combination of the new primitive vectors.

\section*{2. Silicon as a lattice and a basis}
Crystalline silicon has the diamond structure: an fcc Bravais lattice with two identical Si atoms in each primitive cell. Using the conventional cubic lattice constant $a$, one may choose
\[
\vct t_1=\frac a2(0,1,1),\quad
\vct t_2=\frac a2(1,0,1),\quad
\vct t_3=\frac a2(1,1,0),
\qquad
\vct\tau_1=(0,0,0),\quad
\vct\tau_2=\frac a4(1,1,1).
\]
The full set of atomic positions is $\vct R+\vct\tau_\nu$, with $\vct R$ an integer combination of the three translations. The conventional cubic cell contains four fcc lattice points and hence eight Si atoms; the primitive cell contains two.

\section*{3. The stated bcc lattice with a two-atom basis}
\begin{sourceNote}
The source explicitly specifies \emph{two atoms per primitive cell}. This is followed below. It is different from a monatomic bcc crystal, which has one atom per primitive cell and two per conventional cubic cell. The unspecified two-atom basis changes the atomic count, but not the Bravais-cell or Brillouin-zone volumes under the stated bcc assumption.
\end{sourceNote}
\subsection*{(a) Atomic number density}
A conventional bcc cube contains $8(1/8)+1=2$ lattice points. Its primitive volume is thus $\Omega=a^3/2$. With two atoms per primitive cell,
\[
\boxed{n_{\rm atom}=\frac2\Omega=\frac4{a^3}}.
\]
If the intended premise had instead been monatomic bcc, the density would be $2/a^3$; that is not the literal premise of this question.

\subsection*{(b) Wigner--Seitz-cell volume}
A Wigner--Seitz cell contains one Bravais lattice point and has the primitive volume:
\[
\boxed{V_{\rm WS}=\frac{a^3}{2}}.
\]

\subsection*{(c) Brillouin-zone volume}
The Brillouin zone is a reciprocal Wigner--Seitz cell, so
\[
\boxed{V_{\rm BZ}=\frac{(2\pi)^3}{\Omega}=\frac{16\pi^3}{a^3}}.
\]

\subsection*{(d) Number of facets}
The reciprocal lattice of bcc is fcc. In units $2\pi/a$, its shortest nonzero reciprocal vectors are
\[
(\pm1,\pm1,0),\quad(\pm1,0,\pm1),\quad(0,\pm1,\pm1).
\]
There are twelve. Each perpendicular-bisector plane
$\vct k\cdot\vct G=|\vct G|^2/2$ supplies one face of the reciprocal Wigner--Seitz cell. The resulting polyhedron is a rhombic dodecahedron with
\[
\boxed{12\ \text{facets}}.
\]
The fourteen-faced truncated octahedron is the \emph{direct-space} Wigner--Seitz cell of bcc, not its first Brillouin zone.

\section*{4. Reciprocal vectors and planar density}
\subsection*{(a) Reciprocal primitive vectors}
The direct primitive volume is $\Omega=a^3/4$. The cross-product formula gives
\[
\vct g_1=\frac{2\pi}{\Omega}(\vct t_2\times\vct t_3),
\]
with cyclic permutations for the other two vectors. Therefore
\[
\boxed{\vct g_1=\frac{2\pi}{a}(-1,1,1),\quad
\vct g_2=\frac{2\pi}{a}(1,-1,1),\quad
\vct g_3=\frac{2\pi}{a}(1,1,-1)}.
\]
For instance, the first direct vector has dot product $2\pi$ with $\vct g_1$ and zero with $\vct g_2,\vct g_3$, verifying the required normalization.

\subsection*{(b) Density on successive lattice planes}
Let $A_{\parallel}$ be the primitive area of the two-dimensional lattice on one plane, containing one lattice point. Successive occupied planes are separated by $d$. The primitive prism obtained from two in-plane translations and one translation to the next plane has volume $A_{\parallel}d=\Omega$; a lateral shift of the next layer does not affect this volume. Thus
\[
\boxed{\sigma=\frac1{A_{\parallel}}=\frac d\Omega=\frac{4d}{a^3}}.
\]
A second derivation counts points in a large slab: $L/d$ layers of area $A$ contain approximately $\sigma A L/d$ points, which must equal $AL/\Omega$. Boundary corrections vanish in the macroscopic limit.

\section*{5. The Born--von Karman sum}
Define
\[
F(\vct k)=k_xk_yk_z\exp(-k_x^2/\alpha-k_y^2/\beta-k_z^2/\gamma).
\]
The allowed wave vectors of a Born--von Karman periodic box form an inversion-symmetric set: if $\vct k$ is allowed, so is $-\vct k$. Since the Gaussian is even and the product of three components is odd,
\[
F(-\vct k)=-F(\vct k).
\]
Every pair cancels, and the unpaired origin contributes zero. The positive parameters ensure absolute convergence, so pairing the infinite sum is legitimate. Consequently
\[
\boxed{\frac1V\sum_{\vct k}F(\vct k)=0}.
\]
This is an exact discrete cancellation, not merely a continuum approximation. In the thermodynamic limit it also follows by replacing $V^{-1}\sum_{\vct k}$ by $(2\pi)^{-3}\int\dd^3k$; each Cartesian odd Gaussian integral vanishes.

\section*{6. Translating a wavefunction between cells}
Put $\vct T=\vct t_m-\vct t_n$. A point $\vct r$ in the cell centered on $\vct t_m$ corresponds to $\vct r-\vct T$ in the cell centered on $\vct t_n$. Bloch's theorem says
\[
\psi_{\vct k}(\vct r)=\ee^{\ii\vct k\cdot\vct T}
\psi_{\vct k}(\vct r-\vct T).
\]
Since the latter function is known as $f$ in the original cell,
\[
\boxed{\psi_{\vct k}(\vct r)
=\ee^{\ii\vct k\cdot(\vct t_m-\vct t_n)}
f\bigl(\vct r-(\vct t_m-\vct t_n)\bigr)}.
\]
Equivalently, with the same local coordinate $\vct\rho$ in both cells,
$\psi(\vct t_m+\vct\rho)=\ee^{\ii\vct k\cdot\vct T}f(\vct t_n+\vct\rho)$.
Both the phase and the shifted function argument are required when $f$ is expressed in the original global coordinate system.

\section*{7. The anisotropic parabolic band}
\subsection*{(a) Critical point and its type}
The gradient and Hessian are
\[
\nabla_{\vct k}\varepsilon=
\left(\frac{2(k_x-k_0)}\alpha,\frac{2k_y}\beta,\frac{2k_z}\gamma\right),
\qquad H=\operatorname{diag}(2/\alpha,2/\beta,2/\gamma).
\]
All three Hessian eigenvalues are positive. Hence
\[
\boxed{\vct k_c=(k_0,0,0),\qquad \varepsilon_c=\varepsilon_0,
\qquad \text{type }M_0\ \text{(minimum)}}.
\]

\subsection*{(b) Density of states}
For $\Delta=E-\varepsilon_0\ge0$, the states below $E$ occupy an ellipsoid with semiaxes
$\sqrt{\alpha\Delta}$, $\sqrt{\beta\Delta}$, and $\sqrt{\gamma\Delta}$. Its volume is
$(4\pi/3)\sqrt{\alpha\beta\gamma}\,\Delta^{3/2}$. There is one allowed wave vector per volume $(2\pi)^3/V$ in reciprocal space. Thus the integrated state count is
\[
\mathcal N(E)=\frac{g_sV\sqrt{\alpha\beta\gamma}}{6\pi^2}
(E-\varepsilon_0)^{3/2}\Theta(E-\varepsilon_0).
\]
Differentiation gives
\[
\boxed{D(E)=\frac{g_sV\sqrt{\alpha\beta\gamma}}{4\pi^2}
\sqrt{E-\varepsilon_0}\,\Theta(E-\varepsilon_0)}.
\]
Use $g_s=1$ for a spin-resolved band, or $g_s=2$ for a band whose two spin states are counted together. Translation of the ellipsoid by $k_0$ leaves its volume unchanged. If the parabola is only a local band-edge approximation, this expression applies near that edge, before other band features or zone boundaries matter.

\section*{8. Three valence electrons per atom in an fcc metal}
An fcc conventional cell contains four atoms. The stated three valence electrons per atom therefore give
\[
n=\frac NV=\frac{12}{a^3}.
\]
By definition, $4\pi(r_sa_0)^3/3=1/n$, where $a_0$ is the Bohr radius. Hence
\[
\boxed{r_s=\frac{a}{(16\pi)^{1/3}a_0}}.
\]
Counting the two spin states in a filled Fermi sphere gives
$n=k_F^3/(3\pi^2)$, so
\[
\boxed{k_F=\frac{(36\pi^2)^{1/3}}a},\qquad
\boxed{\varepsilon_F=\frac{\hbar^2}{2m_ea^2}(36\pi^2)^{2/3}}.
\]
Finally,
\[
\overline\varepsilon
=\frac{\int_0^{k_F}(\hbar^2k^2/2m_e)k^2\dd k}
{\int_0^{k_F}k^2\dd k}
=\frac{\hbar^2}{2m_e}\frac{k_F^5/5}{k_F^3/3}
=\boxed{\frac35\varepsilon_F}.
\]
This uses the stipulated free-electron model; it does not assume that actual multivalent-metal bands are exactly free-electron dispersions.

\section*{9. Sommerfeld expansion at fixed particle number}
The expansion requires $T\ll T_F$ and a density of states smooth near the chemical potential. The density of states is zero below the bottom of the band, so the printed lower integration limit $-\infty$ is harmless.

\subsection*{(a) Temperature derivative of the chemical potential}
Apply the supplied formula with $\Gamma=D$:
\[
N=\int_{-\infty}^{\mu(T)}D(E)\dd E
+\frac{\pi^2}{6}\kb^2T^2D'(\mu(T))+O(T^4).
\]
At fixed $N$ and volume, differentiation gives
\[
0=D(\mu)\mu'
+\frac{\pi^2}{3}\kb^2T D'(\mu)
+\frac{\pi^2}{6}\kb^2T^2D''(\mu)\mu'+O(T^3),
\]
where $\mu'=\dd\mu/\dd T$. The leading change in $\mu$ is quadratic in $T$, so $\mu'=O(T)$ and the third term is $O(T^3)$. Keeping the requested order,
\[
\boxed{\frac{\dd\mu}{\dd T}
=-\frac{\pi^2}{3}\kb^2T\frac{D'(\mu)}{D(\mu)}+O(T^3)}.
\]
The volume is fixed, so no temperature-dependent volume factor is differentiated.

\subsection*{(b) The three-dimensional free-electron result}
For $E>0$, $D(E)=C\sqrt E$ and $D'/D=1/(2E)$. Therefore
\[
\frac{\dd\mu}{\dd T}=-\frac{\pi^2\kb^2T}{6\mu}+O(T^3).
\]
Multiply by $2\mu$ and integrate from $0$ to $T$:
\[
\mu(T)^2=\mu_0^2-\frac{\pi^2}{6}\kb^2T^2+O(T^4).
\]
Expanding the positive square root, and using $\mu_0=\varepsilon_F=\kb T_F$, yields
\[
\boxed{\mu(T)=\mu_0\left[1-\frac{\pi^2}{12}
\left(\frac T{T_F}\right)^2\right]+O(T^4)}.
\]
Here and below, the remainder notation in dimensional formulas carries the appropriate constant energy and temperature scales.

\subsection*{(c) Internal energy}
Set $\Gamma(E)=E D(E)$. Then $\Gamma'(E)=D(E)+E D'(E)$, giving
\[
\boxed{U(T)=\int_{-\infty}^{\mu(T)}E D(E)\dd E
+\frac{\pi^2}{6}(\kb T)^2
\bigl[D(\mu(T))+\mu(T)D'(\mu(T))\bigr]+O(T^4)}.
\]
The upper limit changes with temperature because the number of electrons is held fixed.

\subsection*{(d) Constant-volume electronic heat capacity}
Differentiate the preceding expression along the fixed-$N$ solution $\mu(T)$:
\begin{align*}
C_V&=\mu D(\mu)\mu'
+\frac{\pi^2}{3}\kb^2T\bigl[D(\mu)+\mu D'(\mu)\bigr]+O(T^3)\\
&=-\frac{\pi^2}{3}\kb^2T\mu D'(\mu)
+\frac{\pi^2}{3}\kb^2T\bigl[D(\mu)+\mu D'(\mu)\bigr]+O(T^3)\\
&=\frac{\pi^2}{3}\kb^2T D(\mu)+O(T^3).
\end{align*}
The two terms involving $\mu D'$ cancel. Since $\mu-\mu_0=O(T^2)$, replacing $D(\mu)$ by $D(\mu_0)$ changes only the cubic correction:
\[
\boxed{C_V(T)=\frac{\pi^2}{3}\kb^2T D(\mu_0)+O(T^3)}.
\]

\subsection*{(e) Why the classical value is strongly suppressed}
Deep occupied states cannot be thermally rearranged freely because of Pauli exclusion. Only a shell of width of order $\kb T$ near $\varepsilon_F$ contributes appreciably. In three dimensions,
\[
D(\varepsilon_F)=\frac{3N}{2\varepsilon_F},\qquad
N_{\rm active}\sim D(\varepsilon_F)\kb T=\frac32N\frac T{T_F}.
\]
Estimating an order-$\kb$ contribution per active state reproduces the printed scale
$C_V\sim(3/2)N\kb(T/T_F)$ and explains why it is much smaller than the classical $(3/2)N\kb$.

\textbf{The printed estimate does not have the exact coefficient.} Substitution of the same $D(\varepsilon_F)$ into part (d) instead gives
\[
\boxed{C_V=\frac{\pi^2}{2}N\kb\frac T{T_F}
+O\!\left(N\kb(T/T_F)^3\right)}.
\]
Thus part (e) is a qualitative active-electron argument. Treating its factor $3/2$ as exact would contradict part (d).

\section*{References and checks}
The lattice, state-counting, and Sommerfeld derivations are displayed explicitly above. For the statistical-mechanical assumptions and low-temperature expansion, see
\href{https://www.damtp.cam.ac.uk/user/tong/statphys/statmechhtml/S3.html}{David Tong, Statistical Physics, Chapter 3}.
The accompanying verification script independently counts finite lattice samples and integrates the finite-temperature Fermi distribution to check the limiting coefficients.
